\section{Methodology}
\label{sec:methodology}
We instantiate the two axes of structure motivated in Section~\ref{sec:research-gaps}
on top of one shared recipe, so that concept space and label space can be
varied independently and every difference we later report is attributable to
structure, not to a different backbone, dataset split, or training budget.

\textbf{Setup.} We work on CUB-200-2011~\cite{wah2011caltech} (200 species,
bounding-box crops). The concepts are the class-level attributes obtained with
the procedure of \citet{koh2020concept}: the noisy instance-level CUB attributes
are aggregated to the class level by \emph{majority voting}---an attribute is
assigned to a class when it is present in the majority of that class's
images---and rare attributes are discarded, leaving a denoised set of $109$
directly supervised concepts. Every variant shares the same feature extractor, a
ConvNeXt-Tiny~\cite{liu2022convnet} backbone trained from scratch, so that any
difference reflects the bottleneck design rather than model capacity. As an upper
bound on accuracy, we also train a no-bottleneck black-box (ConvNeXt-Tiny with a
linear head mapping the image directly to the 200 species). In all concept-based
variants a linear classifier maps the concept vector to the species.

\textbf{Flat CBM.} A whole-image ConvNeXt-Tiny encoder predicts all $109$
concepts at once (sigmoid outputs), without part localisation; the resulting
concept vector feeds the species classifier.

\textbf{Part-conditioned CBM: breaking prediction circularity.} To prevent a
concept from being inferred through whole-image, class-correlated shortcuts, we
localise the visual evidence for each concept \emph{before} predicting it. A
frozen Keypoint R-CNN detector finds 12 anatomical parts (beak, wings, tail,
\dots) together with their visibility, producing one crop per part plus a global
crop of the bird. These crops are encoded by ConvNeXt-Tiny backbones (one shared
across parts, one dedicated to the global crop). Each attribute is then predicted
\emph{only} from the crop of the part it names: for instance, \emph{wing color}
sees only the wing crop, so the concept must rely on local evidence and cannot
exploit species-diagnostic features visible elsewhere in the image. The concept
prediction is therefore decoupled from the class prediction, making the
bottleneck faithful. Parts that are not detected are replaced by a learned
``absent'' embedding. By design, this grounding is expected to cost some
accuracy, because it removes the whole-image shortcuts a flat model would
otherwise use. Figure~\ref{fig:part-crops} shows this process on an actual test
image: keypoint localisation with the resulting crop regions, and the
per-part crops it produces.

\begin{figure*}[t]
  \centering
  \includegraphics[width=0.62\textwidth]{part_crops}
  \caption{Part localisation (a) and the resulting per-part crops (b) on a
  CUB-200-2011 test image.}
  \label{fig:part-crops}
\end{figure*}

\textbf{Coarse-to-fine concepts.} On top of the grounded concepts we add
hierarchy \emph{within} the concept space for color and pattern attribute groups.
Instead of predicting the fine value directly, a small coarse head first
estimates a coarse colour family (e.g.\ warm / cool / neutral / iridescent);
its posterior is then concatenated with the part embedding and fed to a fine
head that predicts the specific value. This hierarchy acts on the concepts, not
on the species output.

\textbf{Training the classifier: independent vs.\ sequential.} The linear
species classifier is trained in two regimes that we compare throughout.
In the \emph{independent} regime it is fit on the ground-truth concept labels;
in the \emph{sequential} regime it is fit on the concept probabilities that the
concept encoder actually outputs, so that it can compensate for the encoder's
systematic errors. Both regimes are evaluated on the real concept predictions,
and the concept encoder itself is identical in both.

\textbf{Label-space hierarchy and mistake severity.} The 200 species form the
leaves of a taxonomic tree $\mathcal{T}$ obtained by mapping each class onto
WordNet (species\,$\rightarrow$\,genus\,$\rightarrow$\,family\,$\rightarrow
\cdots$)~\cite{miller1995wordnet}. For two nodes $a,b$ let $\mathrm{LCA}(a,b)$
denote their lowest common ancestor in $\mathcal{T}$, and let the \emph{height}
$h(v)$ of a node $v$ be the number of edges on the longest downward path from
$v$ to a leaf, so that $h(\cdot)=0$ for leaves. Following
\citet{bertinetto2020making}, the severity of predicting $\hat{y}$ when the true
class is $y$ equals the height of their LCA:
\begin{equation}
  s(y,\hat{y}) \;=\; h\!\left(\mathrm{LCA}(y,\hat{y})\right),
  \label{eq:cost}
\end{equation}
which is $0$ for a correct prediction, $1$ when the two species share a genus,
$2$ when they share only a family, and so on. The mistake severity we report is
the mean over all $N$ test examples:
\begin{equation}
  \mathrm{Sev} \;=\; \frac{1}{N}\sum_{i=1}^{N} s\!\left(y_i,\hat{y}_i\right)
  \;=\; \frac{1}{N}\sum_{i=1}^{N} h\!\left(\mathrm{LCA}(y_i,\hat{y}_i)\right),
  \label{eq:sev}
\end{equation}
so correct predictions contribute nothing and the score captures both how often
the model errs and how far its errors fall in the taxonomy (lower is better).

\textbf{Severity-aware decision (CRM).} The same cost drives the decision rule.
Let $p(c\mid x)$ be the model's softmax posterior over the $C{=}200$ classes.
Instead of the arg\,max, CRM~\cite{karthik2021nocost} returns the class of
minimum \emph{expected} severity:
\begin{equation}
  \hat{y}_{\mathrm{CRM}}
  \;=\; \arg\min_{k\in\{1,\dots,C\}} \sum_{c=1}^{C} s(k,c)\,p(c\mid x),
  \label{eq:crm}
\end{equation}
so that, when the model is uncertain, it favours a class that is taxonomically
close to the likely candidates rather than blindly picking the single most
probable one. Following \citet{karthik2021nocost}'s own formulation, we apply
CRM directly to an already-trained classifier's posterior: it is fully
post-hoc, adds no parameters, and requires no retraining.

\textbf{Metrics.} At the species level we report top-1 accuracy and the mistake
severity of Eq.~\eqref{eq:sev} (lower is better), which measures how
semantically close the errors are. At the concept level we report binary accuracy
and macro-averaged $F_1$ over the $109$ attributes for the concept encoder alone,
to check whether a numerically higher concept score reflects genuinely grounded
evidence or merely a class-correlated shortcut.
